% Don't touch this %%%%%%%%%%%%%%%%%%%%%%%%%%%%%%%%%%%%%%%%%%%
\documentclass[11pt]{article}
\usepackage{fullpage}
\usepackage[left=1in,top=1in,right=1in,bottom=1in,headheight=3ex,headsep=3ex]{geometry}
\usepackage{graphicx}
\usepackage{float}
\usepackage[spanish]{babel}

\newcommand{\blankline}{\quad\pagebreak[2]}
%%%%%%%%%%%%%%%%%%%%%%%%%%%%%%%%%%%%%%%%%%%%%%%%%%%%%%%%%%%%%%

% Modify Course title, instructor name, semester here %%%%%%%%

\title{Micreconom\'ia Aplicada (MAE)}
\author{Jos\'e Manuel Paz y Mi\~no}
\date{Oto\~no 2025}

%%%%%%%%%%%%%%%%%%%%%%%%%%%%%%%%%%%%%%%%%%%%%%%%%%%%%%%%%%%%%%

% Don't touch this %%%%%%%%%%%%%%%%%%%%%%%%%%%%%%%%%%%%%%%%%%%
\usepackage[sc]{mathpazo}
\linespread{1.05} % Palatino needs more leading (space between lines)
\usepackage[T1]{fontenc}
\usepackage[mmddyyyy]{datetime}% http://ctan.org/pkg/datetime
\usepackage{advdate}% http://ctan.org/pkg/advdate
\newdateformat{syldate}{\twodigit{\THEMONTH}/\twodigit{\THEDAY}}
\newsavebox{\MONDAY}\savebox{\MONDAY}{Mon}% Mon
\newcommand{\week}[1]{%
%  \cleardate{mydate}% Clear date
% \newdate{mydate}{\the\day}{\the\month}{\the\year}% Store date
  \paragraph*{\kern-2ex\quad #1, \syldate{\today} - \AdvanceDate[4]\syldate{\today}:}% Set heading  \quad #1
%  \setbox1=\hbox{\shortdayofweekname{\getdateday{mydate}}{\getdatemonth{mydate}}{\getdateyear{mydate}}}%
  \ifdim\wd1=\wd\MONDAY
    \AdvanceDate[7]
  \else
    \AdvanceDate[7]
  \fi%
}
\usepackage{setspace}
\usepackage{multicol}
%\usepackage{indentfirst}
\usepackage{fancyhdr,lastpage}
\usepackage{url}
\pagestyle{fancy}
\usepackage{hyperref}
\usepackage{lastpage}
\usepackage{amsmath}
\usepackage{layout}

\usepackage{natbib}

\lhead{}
\chead{}
%%%%%%%%%%%%%%%%%%%%%%%%%%%%%%%%%%%%%%%%%%%%%%%%%%%%%%%%%%%%%%

% Modify header here %%%%%%%%%%%%%%%%%%%%%%%%%%%%%%%%%%%%%%%%%
\rhead{\footnotesize Microeconom\'ia Aplicada (MAE)}

%%%%%%%%%%%%%%%%%%%%%%%%%%%%%%%%%%%%%%%%%%%%%%%%%%%%%%%%%%%%%%
% Don't touch this %%%%%%%%%%%%%%%%%%%%%%%%%%%%%%%%%%%%%%%%%%%
\lfoot{}
\cfoot{\small \thepage/\pageref*{LastPage}}
\rfoot{}

\usepackage{array, xcolor}
\usepackage{color,hyperref}
\definecolor{clemsonorange}{HTML}{EA6A20}
\hypersetup{colorlinks,breaklinks,linkcolor=clemsonorange,urlcolor=clemsonorange,anchorcolor=clemsonorange,citecolor=black}


\begin{document}

\maketitle

\blankline

\setlength{\parindent}{0pt}.
\begin{tabular*}{.93\textwidth}{@{\extracolsep{\fill}}lr}

%%%%%%%%%%%%%%%%%%%%%%%%%%%%%%%%%%%%%%%%%%%%%%%%%%%%%%%%%%%%%%

% Modify information %%%%%%%%%%%%%%%%%%%%%%%%%%%%%%%%%%%%%%%%%
E-mail: \texttt{jpazm@fen.uchile.cl} & Web: Canvas / Mis Cursos FEN  \\

 Oficina: 1503 \\
 Horas de Oficina: Agendar por correo &  D\'ias y Horas: Mar/Jue 08:00-09:20 \\
 & \\
% Lab Room: ... & Lab Hours: W 3-5pm \\
% &  \\
\hline
\end{tabular*}

\vspace{5 mm}

% First Section %%%%%%%%%%%%%%%%%%%%%%%%%%%%%%%%%%%%%%%%%%%%

\section*{Descripci\'on}

Este es el primer curso de microeconom\'ia del Magister en Econom\'ia Aplicada. El curso tiene 2 prop\'ositos: i. repasar y extender los principales desaarrollos te\'oricos de teor\'ia de juegos y ii. enfatizar aplicaciones pasando por organizaci\'on industrial, econom\'ia politica, salud, contratos, medio ambiente, entre otros. Se espera que el alumno/a no solo pueda identificar los principales modelos te\'oricos de su campo de interes, si no que tambien sea capaz de modificarlos y aumentarlos seg\'un sus requerimientos.

\bigskip

% Second Section %%%%%%%%%%%%%%%%%%%%%%%%%%%%%%%%%%%%%%%%%%%

\section*{Evaluaciones}
2 evaluaciones en clase (solemne y examen final) con un peso de 30\% cada uno en la nota final. 2 tareas para la casa ha realizarse en grupos de 2 con un peso de 10\% cada una en la nota final. Adicionalmente, un proyecto de investigaci\'on (tambi\'en en grupos de 2) con un peso de 20\%. M\'as detalles sobre el proyecto de investigaci\'on se daran en clase.

\section*{Horarios de Oficina}
Coordinar por correo.

\section*{Prerrequisitos}
En general, tener conocimientos s\'olidos de teor\'ia de juegos (lo que en la FEN corresponde al curso de Econom\'ia Aplicada).

\section*{Temas y Referencias}

\begin{enumerate}

  \item T\'opico 0 Introducci\'on.

  \item T\'opico 1 Juegos est\'aticos con informaci\'on completa. Referencias: \cite{osborne1994course} (2,3), \cite{fudenberg1991game} (1,2), \cite{tadelis2013game} (3,5,6), \cite{kreps2023microeconomic} (A9).

  \item T\'opico 2 Juegos din\'amicos con informaci\'on completa. Referencias: \cite{osborne1994course} (6,7), \cite{fudenberg1991game} (3), \cite{tadelis2013game} (7,8,9), \cite{kreps2023microeconomic} (A10-A11).


  \item T\'opico 3 Juegos est\'aticos con informaci\'on incompleta. Referencias: \cite{osborne1994course} (2.6), \cite{fudenberg1991game} (6), \cite{tadelis2013game} (12).


  \item T\'opico 4 Juegos din\'amicos con informaci\'on incompleta. Referencias: \cite{osborne1994course} (11-12), \cite{fudenberg1991game} (8), \cite{tadelis2013game} (15,16,18), \cite{kreps2023microeconomic} (A9).


  \item T\'opico 5 Juegos repetidos. Referencias: \cite{osborne1994course} (8), \cite{fudenberg1991game} (4,5), \cite{tadelis2013game} (10,11), \cite{kreps2023microeconomic} (A15).

  \item T\'opico 6 Otros (reputaci\'on, contratos). Referencias: A determinar.
\end{enumerate}

Otras referencias \'utiles son \cite{mas1995microeconomic}, \cite{mailath2006repeated}, \cite{kartiklecture2012} y \cite{mailath2018modeling}.

\bibliographystyle{ecca}
\bibliography{references_micro_mae}



\end{document}
